\documentclass[11pt]{article}
\usepackage[margin=0.92in]{geometry}
\usepackage{graphicx}
\usepackage{booktabs}
\usepackage{amsmath}
\usepackage{amssymb}
\usepackage{fancyhdr}
\usepackage[hidelinks]{hyperref}
\usepackage{caption}
\usepackage{float}
\usepackage{lastpage}
\usepackage{enumitem}
\usepackage{titlesec}

\titlespacing*{\section}{0pt}{7pt}{2pt}
\captionsetup{font=small,skip=4pt}
\setlength{\intextsep}{6pt}

\pagestyle{fancy}
\fancyhead[L]{}
\fancyhead[R]{\textit{HC Adaptive Retrieval -- Amazon Datasets: Summary}}
\fancyfoot[C]{\thepage/\pageref{LastPage}}

\setlength{\parindent}{0pt}
\setlength{\parskip}{3pt}

\begin{document}

\begin{center}
{\LARGE \textbf{Higher Criticism Adaptive Retrieval on Amazon Queries}} \\[6pt]
{\large From a Controlled Proof-of-Concept to a Calibration Breakthrough}
\end{center}

\vspace{4pt}

% =====================================================================
\section{Motivation: Controlled Environments for HC}
% =====================================================================

Higher Criticism (HC) replaces fixed top-$k$ retrieval by \emph{statistically}
deciding how many documents to return per query: cosine similarities are turned
into p-values against a null distribution, and the HC statistic locates the
threshold where relevant documents stop concentrating. The aim of this effort
was not merely to benchmark HC, but to build \textbf{controlled environments}
in which its behaviour can be \emph{understood}. We wanted datasets that
(i)~emulate a realistic retrieval setting (Amazon product search),
(ii)~have fully known ground-truth structure --- every document's relevance
tier and every query's true $K$ are fixed by construction --- and
(iii)~expose a single, tunable difficulty knob. With these three properties
we can isolate the conditions under which HC works correctly and the conditions
under which it fails, and attribute each failure to a specific mechanism rather
than guess at it. The work follows exactly this axis: a dataset engineered to
\emph{fit} HC's core assumption, then harder datasets that \emph{controllably
violate} it, and finally a diagnostic study that turns the known structure into
a precise account of failure --- and a fix.

% =====================================================================
\section{Proof of Concept: the category-key Dataset}
% =====================================================================

The \texttt{category\_key} dataset is the regime built to satisfy HC's
assumptions. Each of 44 queries asks ``What are the largest cities in
[country]?''; the corpus holds 726 documents (506 relevant city passages,
220 distractors). The number of relevant documents varies by query --- small
($K\!=\!4$--5), medium ($K\!=\!9$--11), and large ($K\!=\!18$--24) buckets ---
so no single fixed $k$ can serve all queries, making adaptivity necessary by
design. Crucially, relevant passages genuinely cluster at high cosine
similarity (BGE-\texttt{base-en-v1.5} embeddings), and a \emph{per-query}
empirical null is near-perfectly calibrated (pooled KS $=0.003$, 44/44 queries
pass). This is the environment in which HC \emph{should} work --- and it does:

\begin{table}[H]
\centering
\begin{tabular}{lccc}
\toprule
\textbf{Method} & \textbf{Recall} & \textbf{Precision} & \textbf{F1} \\
\midrule
Top-10 (best fixed-$k$) & 0.800 & 0.766 & 0.783 \\
Top-20                  & 0.939 & 0.511 & 0.662 \\
\midrule
\textbf{HC (mean $k\!=\!9.7$)} & \textbf{0.852} & \textbf{0.938} & \textbf{0.893} \\
\bottomrule
\end{tabular}
\caption{\texttt{category\_key}: HC beats every fixed-$k$ baseline on F1 by
adapting $k$ per query (true-$K$ correlation $0.835$). This validates the HC
machinery and gives a clean reference point.}
\end{table}

% =====================================================================
\section{Turning the Difficulty Knob: Compound Queries}
% =====================================================================

To stress HC in a \emph{controlled} way, we introduced \textbf{compound
queries} that combine a category filter with a numerical price constraint,
e.g.\ ``Find Electronics products priced above \$150.'' This is a deliberate,
measurable violation of HC's premise: standard text embeddings match category
semantics but \textbf{cannot reason about numerical inequalities}, so hard
negatives (right category, wrong price) sit right next to relevant documents
in cosine space. The difficulty is a single, well-defined knob, and every
document's tier is known by construction.

Two corpora instantiate this regime (OpenAI \texttt{text-embedding-3-small},
FAISS \texttt{IndexFlatIP}): \texttt{amazon\_compound} --- 80{,}000 documents
across three departments, 135 queries, only 4.8\% relevant --- and the larger
\texttt{electronics\_100k} --- 100{,}000 single-department documents, 44
queries, 1.2\% relevant, with 44.5\% \emph{uncapped} hard negatives. To scale,
we also adopted the architectural constraint of a \textbf{single global null}
(one calibrated distribution shared by all queries at inference), replacing
\texttt{category\_key}'s per-query nulls.

\textbf{The failure.} Under the calibrated global null, HC breaks on
\texttt{amazon\_compound}: \textbf{55 of 135 queries produce a negative HC
statistic and return $k\!=\!0$} --- nothing at all. The controlled design is
what lets us ask \emph{which} 55, and \emph{why}.

% =====================================================================
\section{Diagnosing the Failure: Three Mechanisms}
% =====================================================================

Because relevance and true $K$ are known, the 135 queries can be partitioned
exactly and each failure attributed to a mechanism:

\begin{enumerate}[leftmargin=2em]
\item \textbf{Embedding limitation (24 ``irreducible'' queries).} For these the
price signal is simply absent from the cosine: relevant documents are scattered
through the ranks (similarity gap $\approx 0$), below the Donoho--Jin detection
boundary. \emph{No} statistical test can recover them --- the pipeline failed
before HC ran.
\item \textbf{Calibration mismatch (31 ``fixable'' queries).} Here the embedding
\emph{does} separate, but the per-query non-relevant z-score spread
($\sigma_q\!\approx\!1.74$) is much narrower than the global null
($\sigma_{\mathrm{null}}\!\approx\!2.06$). This inflates p-values and trips the
negative-HC gate. The std-ratio mismatch is the \emph{dominant} predictor of HC
failure ($r\!=\!0.83$ with the HC statistic).
\item \textbf{Signal dilution (general).} With a fixed scan window $\gamma$, the
HC win condition $p_{(i)}\!<\!i/N$ tightens linearly as the pool $N$ grows, so a
bigger pool is not a free fix --- it trades calibration gains for dilution
losses.
\end{enumerate}

The key realization: the 31 fixable queries are broken by \emph{miscalibration},
not by missing signal --- and the miscalibration is a direct consequence of the
global-null constraint we adopted for scale. The fix must therefore
\emph{recalibrate per query without abandoning the single global null}.

% =====================================================================
\section{The Breakthrough: \texttt{rescale\_trim}}
% =====================================================================

\texttt{rescale\_trim} is a label-free, per-query transform applied to the
z-scores \emph{before} the global-null lookup. For each query we estimate a
robust noise scale $\hat{\sigma}_q$ as the standard deviation of the pool
z-scores \emph{after dropping the top $\gamma$ ranks} (where relevant documents
concentrate and would otherwise contaminate the estimate), then rescale:
\[
  z'_i \;=\; z_i \cdot \frac{\sigma_{\mathrm{null}}}{\hat{\sigma}_q},
  \qquad
  \hat{\sigma}_q = \mathrm{std}\!\left(z_{(\gamma+1:N)}\right).
\]
This maps every query onto the same scale as the global null, so the empirical
CDF lookup becomes meaningful, and it directly cancels the $r\!=\!0.83$
std-ratio mismatch identified above. The transform uses only observable pool
statistics (no labels) and is applied uniformly to every query --- the null
\emph{itself} stays a single global distribution, so the architectural
constraint is preserved.

\begin{table}[H]
\centering
\begin{tabular}{lrccc}
\toprule
\textbf{Group} & \textbf{n} & \textbf{Top-50} & \textbf{Global null (no rescale)} & \textbf{HC + \texttt{rescale\_trim}} \\
\midrule
positive\_hc & 77 & 0.233 & 0.253 & \textbf{0.251} \\
fixable      & 31 & 0.164 & 0.000 & \textbf{0.157} \\
irreducible  & 24 & 0.014 & 0.000 & 0.012 \\
\bottomrule
\end{tabular}
\caption{Per-group macro-F1 on \texttt{amazon\_compound}. \texttt{rescale\_trim}
\textbf{recovers 30 of 31 fixable queries} ($0.000\!\to\!0.157$), preserves the
77 already-working queries (within $0.002$), and correctly leaves the 24
embedding-limited queries untouched. The failure mode is substantially
eliminated --- without touching the global null and without using any labels.}
\end{table}

% =====================================================================
\section{Final Results and Generalization}
% =====================================================================

\textbf{Best-in-class retrieval.} On \texttt{amazon\_compound},
HC\,+\,\texttt{rescale\_trim} is the \emph{only} method whose F1 strictly beats
every fixed top-$k$ baseline, at an adaptive mean $k\!=\!34$. It also leads a
panel of competing selection rules evaluated on the identical pipeline:

\begin{table}[H]
\centering
\begin{tabular}{lcccc}
\toprule
\textbf{Method} & \textbf{Recall} & \textbf{Precision} & \textbf{F1} & \textbf{Mean $k$} \\
\midrule
Top-20 (fixed)                       & 0.179 & 0.192 & 0.185 & 20.0 \\
Top-50 (fixed)                       & 0.287 & 0.139 & 0.187 & 50.0 \\
\midrule
\textbf{HC + \texttt{rescale\_trim}} & 0.209 & \textbf{0.209} & \textbf{0.209} & 34.0 \\
Berk--Jones + \texttt{rescale\_trim} & 0.238 & 0.176 & 0.202 & 34.8 \\
Adaptive-$k$ (2025)                  & 0.104 & 0.204 & 0.138 & 9.0 \\
Cluster / CAR (2025)                 & 0.598 & 0.060 & 0.110 & 416.5 \\
Benjamini--Hochberg                  & 0.064 & 0.087 & 0.074 & 1.6 \\
Bonferroni                           & 0.014 & 0.049 & 0.022 & 0.3 \\
\bottomrule
\end{tabular}
\caption{Algorithm comparison at the canonical pool $N\!=\!1000$. HC posts the
best F1 and highest precision; Berk--Jones follows closely. The FWER/FDR rules
collapse (flat thresholds cannot find an adaptive cut-off); CAR reaches HC's F1
only when its pool is hand-tuned to $N\!=\!100$, whereas HC needs no pool-size
hyperparameter.}
\end{table}

\textbf{Generalization to scale.} On the larger, single-department
\texttt{electronics\_100k} (100K docs, 44.5\% uncapped hard negatives), HC again
attains the best retrieval F1 ($0.192$ vs.\ $0.185$ top-20, $0.166$ top-50),
$+34.6\%$ recall over top-20 at adaptive $k\!=\!34$. The advantage carries
through the full RAG chain: against top-20, HC lifts \textbf{Token~F1 +15.0\%,
Title~F1 +19.6\%, and Fuzzy-Title~F1 +17.0\%} (\texttt{gpt-5-mini} generation,
\texttt{gpt-5.2} judge), reaching 94--97\% of top-50's generation quality with
32\% fewer documents --- the adaptive-$k$ benefit is not a small-dataset artifact.

\begin{table}[H]
\centering
\begin{tabular}{llccc}
\toprule
\textbf{Dataset} & \textbf{Regime} & \textbf{Docs / Queries} & \textbf{Best fixed-$k$ F1} & \textbf{HC F1} \\
\midrule
\texttt{category\_key}   & clean semantic, var-$K$        & 726 / 44      & 0.783 & \textbf{0.893} \\
\texttt{amazon\_compound}& category+price, 3 depts        & 80K / 135     & 0.192 & \textbf{0.209} \\
\texttt{electronics\_100k}& category+price, 1 dept (100K)  & 100K / 44     & 0.185 & \textbf{0.192} \\
\bottomrule
\end{tabular}
\caption{The journey at a glance. HC is the best method in every regime. The
absolute F1 drop from \texttt{category\_key} to the compound corpora reflects
\emph{dataset difficulty} (embeddings blind to price, $\le$4.8\% relevant), not
a degradation of HC --- which remains the top performer throughout.}
\end{table}

% =====================================================================
\section{Conclusions}
% =====================================================================

The controlled-environment strategy paid off. Engineering datasets with known
structure and a single difficulty knob let us (1)~confirm HC's adaptive-$k$
advantage where its assumptions hold (\texttt{category\_key}), (2)~\emph{decompose}
its failure on hard compound queries into embedding-limited, calibration-mismatched,
and dilution components, and (3)~design \texttt{rescale\_trim} --- a label-free
per-query recalibration that eliminates the dominant (calibration) failure mode
while respecting the single-global-null constraint. The outcome is the best retrieval
F1 of any method tested, generalizing from 726 to 100{,}000 documents and through to
end-to-end RAG quality. The residual ``irreducible'' failures are a property of the
embeddings, not of HC, and point to the natural next step: hybrid retrieval that
combines embeddings for category semantics with structured metadata filtering
for the numeric constraint.

\end{document}
